% Chinese abstract.

\begin{abstractzh}

随着大语言模型长上下文推理窗口不断扩大，自回归推理过程中键值缓存的显存占用不断增加。在真实模型部署场景中，算力硬件的显存逐渐难以满足推理需要。虽然降低缓存位宽能够直接减少显存占用，但因为 Key 侧扰动可能经 logits 与 softmax 改写注意力分布，Value 侧扰动通过加权聚合影响输出表示，现有量化路径常用的张量级重建误差校准并不一定能充分展示量化缓存进入注意力计算后的功能偏移。本文提出一种面向高效推理的 KV Cache 行为对齐量化框架，将注意力分布、聚合输出与任务表现作为共同的校准标准，连接量化参数选择、低比特恢复和量化预算分配三个决策环节。

本文使用校准样本评估全精度参考配置与候选量化配置，记录 attention-distribution KL 与注意力聚合输出扰动两类读数用于筛选量化校准参数，保留下游任务读数作为任务层面的审计依据。然后把校准阶段得到的逐层读数进一步汇总为层级行为敏感度，形成跨模型维度敏感度画像图谱，从而观察不同模型的关键层分布形态。并利用画像图谱在固定 $k$ 预算下比较行为引导分配与位置启发式基线，提出通过 \texttt{AutoK} 根据敏感度排序和覆盖率阈值自动生成模型级 $k$ 候选。构建了一套离线行为读数，同时实现保真校准、低比特行为恢复、跨模型结构判读和预算候选生成，使量化路径与层间预算建议形成连续证据链。

本文的实验覆盖 Qwen2.5、LLaMA-3.1 与 Mistral 三个模型族的六个开源指令模型，结果显示 KV Cache 量化和预算分配都具有明显的模型族、规模和任务依赖性。行为引导校准在保守量化位宽INT8下能够保持稳定保真，进入 \texttt{INT4} 后，对称量化在 Qwen 系列中出现检索阶跃崩塌，K/V 角色诊断表明 Key 侧低比特噪声是直接的失稳触发源，采用角色感知低比特路径 \texttt{INT4-RoleAlign} 可将检索行为恢复到可用区间。跨模型预算实验中自动参数选择 \texttt{AutoK} 效果最好的是模型Mistral-7B，Qwen2.5-3B 对应早层保护有效区间，LLaMA-3.1-8B 呈现中等规模共识区间，Qwen2.5-14B 则表现为高性能簇接近而无稳定最优。\texttt{INT4} 路径可带来约 $73.4\%$ 的 KV Cache 容量压缩，但融合解码的时间收益受 $H_{kv}$、序列长度和后端实现共同影响。
总体而言，本文验证了一条以注意力行为保持为核心、以 K/V 角色差异恢复为低比特路径、以 \texttt{AutoK} 和行为敏感度画像为预算接口的 KV Cache 行为对齐量化框架。

\keywordszh{大语言模型；KV Cache 量化；注意力行为保持；低比特推理；行为引导校准；预算分配}

\end{abstractzh}
